% ========== DESKTOP FRAMEWORK: TAURI 2.x ==========
% Source: Symphony/Content/Tech Stack Guide, Architecture documents
% Tauri 2.x architecture, native API integration, and cross-platform considerations

\lettrine{T}{auri} 2.x serves as Symphony's desktop application framework, providing a modern alternative to Electron that combines the flexibility of web technologies with the performance and security benefits of native applications. This architectural choice reflects Symphony's commitment to delivering a lightweight, secure, and performant development environment.

\subsection*{Architecture Overview}

Tauri's architecture fundamentally differs from traditional web-to-desktop solutions by maintaining a clear separation between the frontend web view and the backend Rust core. This design enables Symphony to leverage web technologies for the user interface while maintaining native performance for core operations.

\begin{figure}[htbp]
    \centering
    \begin{tikzpicture}[node distance=2cm, auto]
        % Frontend Layer
        \node[draw, rectangle, fill=brandPrimary!20, minimum width=8cm, minimum height=1.5cm] (frontend) {
            \textbf{Frontend Layer (WebView)}
        };
        \node[below=0.3cm of frontend.south west, anchor=north west, font=\small] {
            React 19.1.0 + Vite + Tailwind CSS + TypeScript
        };
        
        % IPC Layer
        \node[draw, rectangle, fill=brandAccent!20, minimum width=8cm, minimum height=1cm, below=1.5cm of frontend] (ipc) {
            \textbf{IPC Bridge Layer}
        };
        \node[below=0.2cm of ipc.south west, anchor=north west, font=\small] {
            JSON-RPC + Event System + Type-Safe Bindings
        };
        
        % Backend Layer
        \node[draw, rectangle, fill=brandTertiary!20, minimum width=8cm, minimum height=1.5cm, below=1.5cm of ipc] (backend) {
            \textbf{Rust Backend Core}
        };
        \node[below=0.3cm of backend.south west, anchor=north west, font=\small] {
            Tokio Runtime + Native APIs + File System + AI Integration
        };
        
        % System Layer
        \node[draw, rectangle, fill=brandSecondary!20, minimum width=8cm, minimum height=1cm, below=1.5cm of backend] (system) {
            \textbf{Operating System}
        };
        \node[below=0.2cm of system.south west, anchor=north west, font=\small] {
            Windows + macOS + Linux
        };
        
        % Arrows
        \draw[<->, thick, brandPrimary] (frontend.south) -- (ipc.north);
        \draw[<->, thick, brandAccent] (ipc.south) -- (backend.north);
        \draw[<->, thick, brandTertiary] (backend.south) -- (system.north);
        
        % Labels
        \node[right=0.5cm of frontend.east] {\small UI Rendering};
        \node[right=0.5cm of ipc.east] {\small Communication};
        \node[right=0.5cm of backend.east] {\small Core Logic};
        \node[right=0.5cm of system.east] {\small Native Access};
    \end{tikzpicture}
    \caption{Tauri 2.x Architecture Layers in Symphony}
\end{figure}

The architecture provides several key advantages for Symphony:

\begin{infobox}[title=Tauri Architecture Benefits]
\begin{compactlist}
    \item \textbf{Memory Efficiency}: 10-50x smaller memory footprint compared to Electron
    \item \textbf{Security Model}: Sandboxed frontend with controlled backend access
    \item \textbf{Native Performance}: Rust backend provides near-native execution speed
    \item \textbf{Bundle Size}: Significantly smaller application bundles (10-20MB vs 100-200MB)
\end{compactlist}
\end{infobox}

\subsection*{Native API Integration}

Symphony leverages Tauri's extensive native API integration to provide seamless access to operating system functionality. This integration is crucial for Symphony's role as a development environment that needs deep system access.

\begin{lstlisting}[language=Rust, caption=Native API Integration in Symphony]
// File System Operations
use tauri::{api::file, AppHandle, Manager};
use std::path::PathBuf;

#[tauri::command]
async fn read_project_files(app_handle: AppHandle, project_path: String) -> Result<Vec<FileNode>, String> {
    let path = PathBuf::from(project_path);
    
    // Validate path security
    if !is_safe_path(&path) {
        return Err("Invalid path".to_string());
    }
    
    // Read directory structure
    let files = file::read_dir(&path, true)
        .map_err(|e| format!("Failed to read directory: {}", e))?;
    
    // Convert to Symphony file nodes
    let file_nodes = files.into_iter()
        .map(|entry| FileNode::from_path_entry(entry))
        .collect();
    
    Ok(file_nodes)
}

// Process Management
#[tauri::command]
async fn spawn_ai_process(command: String, args: Vec<String>) -> Result<ProcessId, String> {
    use std::process::Command;
    
    let mut cmd = Command::new(&command);
    cmd.args(&args);
    
    // Configure process environment
    cmd.env("SYMPHONY_MODE", "ai_worker");
    cmd.env("RUST_LOG", "info");
    
    let child = cmd.spawn()
        .map_err(|e| format!("Failed to spawn process: {}", e))?;
    
    Ok(ProcessId(child.id()))
}

// System Information
#[tauri::command]
async fn get_system_info() -> SystemInfo {
    SystemInfo {
        os: std::env::consts::OS.to_string(),
        arch: std::env::consts::ARCH.to_string(),
        cpu_count: num_cpus::get(),
        memory_total: sys_info::mem_info().unwrap().total,
        available_space: get_available_disk_space(),
    }
}
\end{lstlisting}

The native API integration extends to platform-specific functionality:

\begin{table}[h]
\centering
\caption{Platform-Specific Native API Usage}
\begin{tabular}{@{}lll@{}}
\toprule
\textbf{Feature} & \textbf{Windows} & \textbf{macOS/Linux} \\
\midrule
File Associations & Registry Integration & Desktop Files \\
System Tray & Windows API & GTK/Cocoa \\
Notifications & WinRT Toast & libnotify/NSUserNotification \\
Window Management & Win32 API & X11/Wayland/Cocoa \\
Shell Integration & PowerShell/CMD & Bash/Zsh/Fish \\
\bottomrule
\end{tabular}
\end{table}

\subsection*{WebView Management}

Symphony's WebView management strategy optimizes for both performance and security. The WebView serves as the rendering engine for Symphony's React-based interface while maintaining strict security boundaries.

\begin{lstlisting}[language=Rust, caption=WebView Configuration and Management]
// WebView Configuration
use tauri::{WebviewUrl, WebviewWindowBuilder, Manager};

fn create_main_window(app: &AppHandle) -> tauri::Result<()> {
    let window = WebviewWindowBuilder::new(
        app,
        "main",
        WebviewUrl::App("index.html".into())
    )
    .title("Symphony - AI-First Development Environment")
    .inner_size(1400.0, 900.0)
    .min_inner_size(800.0, 600.0)
    .resizable(true)
    .decorations(true)
    .transparent(false)
    .shadow(true)
    .center()
    .build()?;
    
    // Configure WebView security
    window.with_webview(|webview| {
        #[cfg(target_os = "linux")]
        {
            use webkit2gtk::WebViewExt;
            webview.inner().set_zoom_level(1.0);
        }
        
        #[cfg(target_os = "windows")]
        {
            use webview2_com::Microsoft::Web::WebView2::Win32::ICoreWebView2Settings;
            // Configure security settings
        }
    })?;
    
    Ok(())
}

// Multi-Window Management
#[tauri::command]
async fn create_ai_panel_window(app_handle: AppHandle) -> Result<(), String> {
    let ai_window = WebviewWindowBuilder::new(
        &app_handle,
        "ai_panel",
        WebviewUrl::App("ai-panel.html".into())
    )
    .title("Symphony AI Panel")
    .inner_size(400.0, 600.0)
    .resizable(true)
    .always_on_top(true)
    .build()
    .map_err(|e| format!("Failed to create AI panel: {}", e))?;
    
    Ok(())
}
\end{lstlisting}

\subsection*{Cross-Platform Considerations}

Symphony's cross-platform strategy addresses the unique requirements and constraints of each target operating system while maintaining a consistent user experience.

\begin{successbox}
\textbf{Cross-Platform Strategy:} Symphony targets Windows 10+, macOS 10.15+, and Linux distributions with modern WebKit support. The application maintains 95\% code sharing across platforms while optimizing for platform-specific user experience patterns.
\end{successbox}

Platform-specific optimizations include:

\begin{lstlisting}[language=Rust, caption=Platform-Specific Optimizations]
// Platform-Specific Window Behavior
#[cfg(target_os = "windows")]
fn configure_windows_specific(window: &WebviewWindow) -> tauri::Result<()> {
    use windows::Win32::UI::WindowsAndMessaging::*;
    
    // Enable Windows 11 rounded corners
    window.set_decorations(true)?;
    
    // Configure Windows-specific shortcuts
    window.register_global_shortcut("Ctrl+Shift+P", || {
        // Open command palette
    })?;
    
    Ok(())
}

#[cfg(target_os = "macos")]
fn configure_macos_specific(window: &WebviewWindow) -> tauri::Result<()> {
    use tauri::TitleBarStyle;
    
    // Use macOS-style title bar
    window.set_title_bar_style(TitleBarStyle::Overlay)?;
    
    // Configure macOS-specific shortcuts
    window.register_global_shortcut("Cmd+Shift+P", || {
        // Open command palette
    })?;
    
    Ok(())
}

#[cfg(target_os = "linux")]
fn configure_linux_specific(window: &WebviewWindow) -> tauri::Result<()> {
    // Configure for various Linux desktop environments
    if let Ok(desktop) = std::env::var("XDG_CURRENT_DESKTOP") {
        match desktop.as_str() {
            "GNOME" => configure_gnome_integration(window)?,
            "KDE" => configure_kde_integration(window)?,
            _ => configure_generic_linux(window)?,
        }
    }
    
    Ok(())
}
\end{lstlisting}

\subsection*{Security Model}

Tauri's security model provides multiple layers of protection that are essential for Symphony's role as a development environment handling sensitive code and AI interactions.

\begin{alertbox}
\textbf{Security Architecture:} Symphony implements a zero-trust security model where the frontend has no direct access to system resources. All system operations must be explicitly exposed through Tauri commands with proper validation and authorization.
\end{alertbox}

The security implementation includes:

\begin{lstlisting}[language=Rust, caption=Security Model Implementation]
// Command Authorization
use tauri::{command, AppHandle, State};
use std::collections::HashMap;

#[derive(Default)]
struct SecurityContext {
    permissions: HashMap<String, Vec<String>>,
    active_sessions: HashMap<String, SessionInfo>,
}

#[tauri::command]
async fn secure_file_operation(
    app_handle: AppHandle,
    security_context: State<'_, SecurityContext>,
    operation: String,
    path: String,
    session_id: String,
) -> Result<String, String> {
    // Validate session
    let session = security_context.active_sessions
        .get(&session_id)
        .ok_or("Invalid session")?;
    
    // Check permissions
    let user_permissions = security_context.permissions
        .get(&session.user_id)
        .ok_or("No permissions found")?;
    
    if !user_permissions.contains(&operation) {
        return Err("Operation not permitted".to_string());
    }
    
    // Validate path
    let canonical_path = std::fs::canonicalize(&path)
        .map_err(|_| "Invalid path")?;
    
    if !is_within_project_bounds(&canonical_path, &session.project_root) {
        return Err("Path outside project bounds".to_string());
    }
    
    // Execute operation with audit logging
    audit_log(&session.user_id, &operation, &path);
    execute_file_operation(&operation, &canonical_path)
}

// Content Security Policy
fn configure_csp() -> String {
    "default-src 'self'; \
     script-src 'self' 'unsafe-inline'; \
     style-src 'self' 'unsafe-inline'; \
     img-src 'self' data: https:; \
     connect-src 'self' ws: wss:; \
     font-src 'self'; \
     object-src 'none'; \
     base-uri 'self'".to_string()
}
\end{lstlisting}

\subsection*{Performance Optimization}

Symphony's Tauri implementation includes several performance optimizations that ensure responsive user experience even during intensive AI processing operations.

\begin{table}[h]
\centering
\caption{Performance Metrics: Symphony vs Electron-based IDEs}
\begin{tabular}{@{}lll@{}}
\toprule
\textbf{Metric} & \textbf{Symphony (Tauri)} & \textbf{Electron IDE} \\
\midrule
Memory Usage (Idle) & 45-60 MB & 200-400 MB \\
Startup Time & 0.8-1.2 seconds & 3-8 seconds \\
Bundle Size & 15-25 MB & 150-300 MB \\
CPU Usage (Idle) & <1\% & 2-5\% \\
\bottomrule
\end{tabular}
\end{table}

Performance optimizations include:

\begin{lstlisting}[language=Rust, caption=Performance Optimization Strategies]
// Async Task Management
use tokio::sync::{mpsc, RwLock};
use std::sync::Arc;

#[derive(Clone)]
struct PerformanceManager {
    task_queue: Arc<RwLock<VecDeque<Task>>>,
    active_tasks: Arc<RwLock<HashMap<TaskId, TaskHandle>>>,
    metrics: Arc<RwLock<PerformanceMetrics>>,
}

impl PerformanceManager {
    async fn execute_with_priority(&self, task: Task) -> Result<TaskResult, TaskError> {
        // Check system resources
        let system_load = self.get_system_load().await;
        if system_load > 0.8 {
            // Defer non-critical tasks
            if task.priority < Priority::High {
                self.defer_task(task).await?;
                return Ok(TaskResult::Deferred);
            }
        }
        
        // Execute with resource monitoring
        let start_time = Instant::now();
        let result = self.execute_task_monitored(task).await?;
        let duration = start_time.elapsed();
        
        // Update performance metrics
        self.update_metrics(duration, result.resource_usage).await;
        
        Ok(result)
    }
    
    async fn optimize_memory_usage(&self) {
        // Garbage collect unused WebView resources
        if let Some(window) = self.get_main_window() {
            window.eval("window.gc && window.gc()").ok();
        }
        
        // Clean up completed tasks
        let mut active_tasks = self.active_tasks.write().await;
        active_tasks.retain(|_, handle| !handle.is_finished());
        
        // Trigger Rust garbage collection if needed
        if self.should_trigger_gc().await {
            // Force garbage collection in Rust (limited options)
            drop(Vec::<u8>::with_capacity(1024 * 1024));
        }
    }
}
\end{lstlisting}

The Tauri 2.x framework provides Symphony with a robust, secure, and performant foundation for desktop application development. By leveraging Rust's performance characteristics and Tauri's security model, Symphony delivers a development environment that is both powerful and resource-efficient, setting a new standard for AI-integrated development tools.